\section*{Envoi}
\addcontentsline{toc}{section}{Envoi}
\label{sec:envoi}

\noindent
This paper has treated a shared flourishing as a cycle to be sustained, a fabric to be tended, and a turning to be appraised in a value that no measure reaches. It has spoken of recyclability and skew, of belief and bifurcation, of habit and household, because a practice requires what is operable, and it has conceded throughout that these are proxies for a good that is lived before it is assessed and that exceeds every assessment of it.

What the assessments leave out is the whole of the matter, and it returns at the close as it returned in the prior paper: a meal shared across years that no satisfaction scale records the worth of, a way of keeping a room that is the form a love has taken, a silence between two people that is full rather than empty and that no instrument can read. These are the matter of the field, and they are the observations through which, beneath all reflection, each appraises continually how things stand between them. The cycle of a shared life turns in them, and turns well when it turns justly, returning to each what each has given, foreclosing the unfolding of neither, and rising, turn upon turn, into more than the two were.

A good so cultivated is not secured and then possessed. It is generated in the turning, and a relation flourishes by continuing to turn, justly, and ceases when the turning ceases. There is no state of it to be reached and kept; there is only the tending, the appraising, the reflecting, and the turning again, across the seasons of a shared life, from its beginning, through whatever it becomes, into its later years.

\begin{flushright}
\itshape\color{rose}
For the forest girl,\\
that our turning, turned together, may rise;\\
that what returns to her be no less than she gives;\\
and that the good of our shared life be generated, and generated again,\\
in every turning of its years.
\end{flushright}

\bigskip
\begin{center}
\itshape\color{rose}
\zh{愿我如星君如月，夜夜流光相皎洁。}
\end{center}
